\subsection{Loss Function}
To adapt the Segment Anything Model (SAM) to the unique characteristics of the ClimateNet dataset, a systematic hyperparameter optimization was conducted. The primary goal of this search was to identify the optimal balance between pixel-wise classification accuracy and regional mask overlap, specifically addressing the extreme class imbalance where background pixels constitute approximately 93.9\% of the data.

\subsubsection{Search Strategy and Configuration
}\label{searchstrat}  

We implemented a random search strategy over 40 iterations to explore the high-dimensional space of the hybrid loss function. To ensure computational efficiency during this search phase, the models were trained on a representative subset consisting of 20\% of the training data. Validation was performed using the ViT-B (Vision Transformer Base) version of the SAM image encoder.

This methodology is grounded in several practical considerations:

\begin{itemize}
\item Efficiency: Training on a 20\% subset allows for a broader exploration of the parameter space within a fixed compute budget.
\item Generalizability: Because ClimateNet consists of global atmospheric rasters with recurring geophysical patterns, a 20\% subset provides sufficient structural diversity—including circular TC structures and ribbon-like AR corridors—to identify robust hyperparameter trends that generalize to the full dataset.
\item Architecture Consistency: Utilizing the ViT-B backbone ensures that the identified hyperparameters are optimized for the specific feature resolution (64 $\times$ 64 feature grid) and embedding space (256 dimensions) used in the final adaptation phase.
\end{itemize}

\subsubsection{Search Space}

The search space was specifically designed to prioritize Recall $\beta$ over Precision $\alpha$ for TCs, as missing a rare storm event is considered a more significant failure in a meteorological context than a slight over-segmentation of its boundary.

\begin{table}[h]
\centering
\begin{tblr}{
  colspec = {llXX}, % X column automatically calculates width and wraps text
  hlines,           % Adds horizontal lines
}
Parameter & Feature & Search Values & Rationale \\ 
Tversky ($\alpha, \beta$) & TC & (0.2, 0.8), (0.1, 0.9), (0.3, 0.7) & Prioritizes recall for rare circular structures. \\ 
Tversky ($\alpha, \beta$) & AR & (0.3, 0.7), (0.4, 0.6), (0.5, 0.5) & Balanced approach for common ribbon forms. \\ 
Focal ($\alpha, \gamma$) & TC & (0.99, 2.0), (0.99, 5.0), (0.995, 8.0) & Combats the 257:1 background-to-TC ratio. \\ 
Focal ($\alpha, \gamma$) & AR & (0.80, 2.0), (0.85, 5.0), (0.90, 2.0) & Addresses the 17:1 background-to-AR ratio. \\ 
\end{tblr}
\caption{Hyperparameter search space for SAM adaptation on ClimateNet.}
\end{table}

% ### Implementation and Stability

During the adaptation phase, the loss weights for the Tversky and Focal components were held constant ($w=1.0$) to ensure gradient stability while the model learned to interpret the 16 multi-parameter atmospheric grids. The search utilized a fixed budget of 50 epochs per iteration. By varying the focusing parameter $\gamma$ up to 8.0, we forced the model to significantly down-weight "easy" background pixels and concentrate on the ambiguous boundaries of the atmospheric features.